\documentclass{article}
\begin{document}
\title{Evidence-Grounded Study of Transparent sentiment baselines across Amazon IMDb and Yelp reviews}
\maketitle

\section{Abstract}
We study Transparent sentiment baselines across Amazon IMDb and Yelp reviews with a local, auditable experiment pipeline. The best kept trial was `ablation-unigram-alpha05-seed1` with primary metric 0.8403343304493236.

\section{Introduction}
The goal is to convert a research idea into a paper candidate without losing provenance.
This draft only reports evidence that exists in the run artifacts.

\section{Related Work}
The screened literature includes Domain adaptation for large-scale sentiment classification: A deep learning approach \cite{glorot2011domainb905a703}, Lexicon-Based Methods for Sentiment Analysis \cite{taboada2011lexiconbasedc1eaa880}, Learning Sentiment-Specific Word Embedding for Twitter Sentiment Classification \cite{tang2014learning98a135fa}, Aspect Level Sentiment Classification with Deep Memory Network \cite{tang2016aspecte780e266}, Deep learning for sentiment analysis: A survey \cite{zhang2018deepc5431074}, End-to-End Adversarial Memory Network for Cross-domain Sentiment Classification \cite{li2017endtoendaadc69e5}, Wasserstein Distance Guided Representation Learning for Domain Adaptation \cite{shen2018wassersteinb4d359ec}, A survey on sentiment analysis methods, applications, and challenges \cite{wankhade2022surveybd7af33b}, SemEval-2014 Task 4: Aspect Based Sentiment Analysis \cite{pontiki2014semeval2014ddaa69ec}, Improving sentiment analysis via sentence type classification using BiLSTM-CRF and CNN \cite{chen2016improving0fec043c}, Aspect-based Sentiment Classification with Aspect-specific Graph Convolutional Networks \cite{zhang2019aspectbased9a079057}, Co-training for cross-lingual sentiment classification \cite{wan2009cotrainingf7cb4166}, Neural Sentiment Classification with User and Product Attention \cite{chen2016neural0df7d563}, Multi-domain sentiment classification \cite{li2008multidomain127a52f4}, A survey of transfer learning \cite{weiss2016surveydb7dbd8b}, Hierarchical Attention Transfer Network for Cross-Domain Sentiment Classification \cite{li2018hierarchical7d97e416}, Learning Semantic Representations of Users and Products for Document Level Sentiment Classification \cite{tang2015learninge778b40f}, A review on sentiment analysis from social media platforms \cite{rodrguezibez2023review357d59da}, Systematic reviews in sentiment analysis: a tertiary study \cite{ligthart2021systematic61a34904}, Sentiment Domain Adaptation with Multiple Sources \cite{wu2016sentiment3761d7c5}, Interactive Attention Transfer Network for Cross-Domain Sentiment Classification \cite{zhang2019interactive2aa8ed3c}, Adversarial and Domain-Aware BERT for Cross-Domain Sentiment Analysis \cite{du2020adversarial0e9baf5d}, Domain Adaptation in Sentiment Classification \cite{uribe2010domain6fd982cc}, Automatic construction of domain-specific sentiment lexicon for unsupervised domain adaptation and sentiment classification \cite{beigi2021automaticef9b17f2}, A review of domain adaptation for opinion detection and sentiment classification \cite{kasthuriarachchy2012review1bc25f5f}, Exploiting Out-of-Domain Datasets and Visual Representations for Image Sentiment Classification \cite{pournaras2021exploitingd282c12a}, T-LBERT with Domain Adaptation for Cross-Domain Sentiment Classification \cite{cao2023tlbert0f143b29}, Attention Based GRU Network for Domain Adaptation in Sentiment Classification \cite{pu2019attention5327674b}, Attention Based GRU Network for Domain Adaptation in Sentiment Classification \cite{pu2019attentione6b21266}, Multi-Domain Aspect-Based Sentiment Analysis Using Domain Adaptation Techniques \cite{toms2025multidomain33316bb2}, An Improved Spectral Feature Alignment for Domain Adaptation in Sentiment Classification \cite{huang2019improved2d2c1374}, Multi-domain adaptation for sentiment classification: Using multiple classifier combining methods \cite{li2008multidomaina820adf0}, A Multi-domain Adaptation for sentiment classification algorithm based on Class Distribution \cite{hu2012multidomainbcc75de0}, An Improved Spectral Feature Alignment for Domain Adaptation in Sentiment Classification \cite{huang2019improveda508c460}, Multi-source domain adaptation with joint learning for cross-domain sentiment classification \cite{zhao2020multisource7b0cdd65}, Domain adaptation with a shrinkable discrepancy strategy for cross-domain sentiment classification \cite{fu2022domain8bfe3f26}, Transfer Learning for Cross-Domain Sentiment Classification via Domain Adaptation \cite{singh2025transfer1f0335ee}, FIT BUT at SemEval-2023 Task 12: Sentiment Without Borders - Multilingual Domain Adaptation for Low-Resource Sentiment Classification \cite{aparovich2023fitc3d3d015}, Adaptation and Use of Subjectivity Lexicons for Domain Dependent Sentiment Classification \cite{dehkharghani2012adaptation0425a16e}, Feature Ensemble Plus Sample Selection: Domain Adaptation for Sentiment Classification \cite{xia2013feature0e6c6b6d}, Datasets for Medical Sentiment Analysis \cite{denecke2023datasets03b63e0c}, Deep Learning Based vs. Markov Chain Based Text Generation for Cross Domain Adaptation for Sentiment Classification \cite{abdelwahab2018deep9757ef1c}, A cross-domain adaptation method for sentiment classification using probabilistic latent analysis \cite{gao2011crossdomainf0de5e2c}, Hyperparameter tuning in convolutional neural networks for domain adaptation in sentiment classification (HTCNN-DASC) \cite{krishnakumari2019hyperparameter922e22d8}, Markov Chain based Method for In-Domain and Cross-Domain Sentiment Classification \cite{domeniconi2015markov6c95be6c}, Multi-Source Domain Adaptation for Visual Sentiment Classification \cite{lin2020multisourcea7234add}, Contrastive transformer based domain adaptation for multi-source cross-domain sentiment classification \cite{fu2022contrastivee139263a}, Cross Domain Sentiment Classification Techniques: A Review \cite{kadli2019cross469bc0a0}, Sentiment Classification Using Convolutional Neural Networks \cite{kim2019sentiment615c3ca2}, Sentiment Analysis in the Era of Large Language Models: A Reality Check \cite{zhang2024sentimentc6eb6828}, Sentiment analysis: A survey on design framework, applications and future scopes \cite{bordoloi2023sentimentdb283d12}, Improving Domain-Adapted Sentiment Classification by Deep Adversarial Mutual Learning \cite{xue2020improvingdba2f6c5}, Cross-lingual sentiment classification in low-resource Bengali language \cite{sazzed2020crosslingual181c8ce3}, Sentiment Classification Using Negative and Intensive Sentiment Supplement Information \cite{chen2019sentimente5e5b9d2}, Weight Poisoning Attacks on Pretrained Models \cite{kurita2020weight048d91fe}, Improving Document-Level Sentiment Classification Using Importance of Sentences \cite{choi2020improving432d296c}, "Challenges and future in deep learning for sentiment analysis: a comprehensive review and a proposed novel hybrid approach" \cite{islam2024challenges02fd95a5}, Halal Products on Twitter: Data Extraction and Sentiment Analysis Using Stack of Deep Learning Algorithms \cite{feizollah2019halalf402d6f1}, Exploring transformer models for sentiment classification: A comparison of <scp>BERT</scp>, <scp>RoBERTa</scp>, <scp>ALBERT</scp>, <scp>DistilBERT</scp>, and <scp>XLNet</scp> \cite{areshey2024exploring4f38bf39}, Sentiment Analysis in the Era of Large Language Models: A Reality Check \cite{zhang2023sentimentb4b66e57}, A Novel Sentiment Analysis for Amazon Data with TSA based Feature Selection \cite{daniel2021novel127fb151}, Analysis of Deep Learning Model Combinations and Tokenization Approaches in Sentiment Classification \cite{erkan2023analysis6633bdc1}, Beneath the Tip of the Iceberg: Current Challenges and New Directions in Sentiment Analysis Research \cite{poria2020beneath22cf40d2}, Exploring the Efficacy of Automatically Generated Counterfactuals for Sentiment Analysis \cite{yang2021exploring1e053e64}, Evolving techniques in sentiment analysis: a comprehensive review \cite{kumar2025evolving5a044b19}, Comparative Analysis of Machine Learning Algorithms for Sentiment Classification in Amazon Reviews \cite{yu2024comparative4e50d20b}, Generalizing sentiment analysis: a review of progress, challenges, and emerging directions \cite{alahmadi2025generalizing369ede0a}, UniGen: Universal Domain Generalization for Sentiment Classification via Zero-shot Dataset Generation \cite{choi2024unigen4fb2ffdf}, Spider Taylor-ChOA: Optimized Deep Learning Based Sentiment Classification for Review Rating Prediction \cite{banbhrani2022spiderd7fe2b64}, Predicting Discourse Structure using Distant Supervision from Sentiment \cite{huber2019predictingfa36d1ac}, Leveraging Lexicon-Based and Sentiment Analysis Techniques for Online Reputation Generation \cite{boumhidi2021leveraging6c43877c}, Polarity of Yelp Reviews: A BERT–LSTM Comparative Study \cite{belaroussi2025polarity26a67026}, Supplemental Information 2: The text based review dataset taken from Amazon, Imdb, and Yelp websites \cite{paperndsupplementald5b276ae}, Supplemental Information 1: The YELP-5 dataset is a sentiment classification task dataset used for sentiment analysis. It comprises user reviews sourced from the Yelp website \cite{paperndsupplemental8b6472a3}, Figure 4: The mean index for the sentiment classification on IMDb dataset. \cite{paperndfigure92a99ba4}, Enhancing Sentiment Analysis with Fine-Tuned DeBERTa in a Cross-Domain Study on IMDb and Amazon Polarity Datasets \cite{cao2025enhancing52f2345f}, Supplemental Information 6: The IMDB dataset is a sentiment classification task dataset used for sentiment analysis. It consists of sentences extracted from movie reviews \cite{paperndsupplementale2a50a6f}, Recurrent Convolutional Neural Network with Attention for Twitter and Yelp Sentiment Classification \cite{wen2018recurrent306cb42f}, Table 3: Extended binary metrics on IMDb and Amazon dataset. \cite{paperndtableaaf66ed6}, Figure 6: The mean index for the sentiment classification on Amazon dataset. \cite{paperndfigure967211c9}, Musiker-Wiki von IMDb und Amazon \cite{gehring2008musikerwikice99da7d}, Musiker-Wiki von IMDb und Amazon \cite{gehring2008musikerwiki92fe96c0}, SENTIMENT FEATURES FOR YELP NOT-RECOMMENDED ONLINE REVIEWS STUDY \cite{lindsentiment6c9bc9d6}, IMDb Movie Data Classification using Voting Classifier for Sentiment Analysis \cite{kaushik2022imdb7b5a4856}, Empirical Study on the Effectiveness of DistilBERT Fine-tuning on IMDb Sentiment Classification Outperforming CNN/LSTM \cite{yuan2026empirical5a975446}, Sentiment Analysis on IMDB Movie Reviews using LSTM \cite{choudhary2025sentimenta52cc157}, Modelling and control of an intelligent communication system for IMDB sentiment classification \cite{yang2025modelling8e1aaa88}, An Empirical Comparison of BERT and Lightweight Variants for IMDb Sentiment Classification \cite{wu2026empirical56d5f7df}, \&lt;span\&gt;Deep Learning-Based Sentiment Analysis: Enhancing IMDb Review Classification with LSTM Models\&lt;/span\&gt; \cite{kalla2025spanc8735c67}, Sentiment Analysis of IMDb Movie Reviews : A comparative study on Performance of Hyperparameter-tuned Classification Algorithms \cite{ghosh2022sentiment2060df3c}, A Lexicon-based Approach for Sentiment Classification of Amazon Books Reviews in Italian Language \cite{chiavetta2016lexiconbased0860c873}, Buy It - Don’t Buy It: Sentiment Classification on Amazon Reviews Using Sentence Polarity Shift \cite{orimaye2012buyd50ada57}, Leveraging Yelp Data for Business Strategy Optimization through Sentiment Analysis and Predictive Modeling \cite{korada2025leveraging6ec8db92}, Deep Learning-Based Sentiment Analysis: Enhancing IMDb Review Classification with LSTM Models \cite{varadharajan2025deep223eb57d}, Sentiment Analysis Research Based on IMDb Movie Reviews \cite{x2026sentiment260218c3}, On Deep Learning in Cross-Domain Sentiment Classification \cite{domeniconi2017deep593e4546}, Binary Sentiment Classification on the IMDB Movie Review Dataset Using an Ensemble Approach \cite{dongre2025binary665b372b}, Yelp Consumption Sentiment and Asset Pricing \cite{jacoby2024yelpa2b73b30}, Cross domain sentiment classification using enhanced sentiment sensitive thesaurus (ESST) \cite{sanju2013cross26b4ce16}, Sentiment Analysis on IMDB Movie Reviews using VADER with Lexical Affinity and Semantic Sentiment Expansion \cite{dsouza2023sentiment9a66f432}, Sentiment Analysis in Social Media Data for Depression Detection Using Artificial Intelligence: A Review \cite{babu2021sentiment5b478c38}, Sentiment Analysis With Ensemble Hybrid Deep Learning Model \cite{tan2022sentiment4832d81e}, An Enhanced Hybrid Feature Selection Technique Using Term Frequency-Inverse Document Frequency and Support Vector Machine-Recursive Feature Elimination for Sentiment Classification \cite{nafis2021enhanced8d424bec}, A Survey of Sentiment Analysis Based on Transfer Learning \cite{liu2019survey6f8551f6}, Impact of convolutional neural network and FastText embedding on text classification \cite{umer2022impact43feb507}, A Hybrid Framework for Sentiment Analysis Using Genetic Algorithm Based Feature Reduction \cite{iqbal2019hybridf5ede317}, Sentiment Analysis of Consumer Reviews Using Deep Learning \cite{iqbal2022sentimentd18b9ffb}, Deep Learning-Based Sentiment Classification: A Comparative Survey \cite{mabrouk2020deepe59e339e}, Detection of Fake News Text Classification on COVID-19 Using Deep Learning Approaches \cite{bangyal2021detection2da80cf9}, Sentiment Analysis using various Machine Learning and Deep Learning Techniques \cite{umarani2021sentiment5263777e}, Document Level Sentiment Analysis: A survey \cite{behdenna2018document0f28a934}, Transformer based Contextual Model for Sentiment Analysis of Customer Reviews: A Fine-tuned BERT \cite{durairaj2021transformer24fb64c1}, Sentiment Analysis on Social Media Reviews Datasets with Deep Learning Approach \cite{baarslan2021sentimentfa026eca}, A Combined Approach of Sentimental Analysis Using Machine Learning Techniques \cite{gupta2023combined9632e9b1}, A Domain-Independent Classification Model for Sentiment Analysis Using Neural Models \cite{jnoub2020domainindependentbc6a4e2b}, BERT: A Review of Applications in Sentiment Analysis \cite{sayeed2023bert678d1641}, Efficient Framework for Sentiment Classification Using Apriori Based Feature Reduction \cite{jain2018efficientdf1b7bd2}, Performance Study of N-grams in the Analysis of Sentiments \cite{ojo2021performanceb18b3098}, A Comparative Study of Different Classification Techniques for Sentiment Analysis \cite{ghosh2020comparative0d185490}, Twenty Years of Machine-Learning-Based Text Classification: A Systematic Review \cite{palanivinayagam2023twentyf98b0ff7}, Exploring Hybrid Approaches for Sentiment Classification: A Comparative Study of LSTM, Naive Bayes, and Bayesian Network on IMDB Reviews \cite{reddy2025exploringf74fbe8e}, IMDb Sentiment Analysis Using Naive Bayes :Multinomial and Bernoulli Models with LaplaceSmoothing and 5-Fold Cross Validation \cite{mishra2025imdb53a34975}, Sentiment Analysis of IMDb Movie Reviews Using SVM and Naive Bayes Classifier \cite{subedi2025sentiment4d1e1933}, Implementation of Sentiment Analysis Movie Review based on IMDB with Naive Bayes Using Information Gain on Feature Selection \cite{mahyarani2021implementation92725ed4}, Opinion Classification on IMDb Reviews Using Naïve Bayes Algorithm \cite{putri2025opinion5a9abee2}, SVM and Naive Bayes for Sentiment Analysis in Arabic \cite{mohamed2022svm467a111d}, The Comparison of Sentiment Analysis of Moon Knight Movie Reviews between Multinomial Naive Bayes and Support Vector Machine \cite{paperndcomparisoneeaa2824}, Sentiment Analysis of Public Reviews of the Brompit Application Using Naive Bayes and SVM Algorithms \cite{affandi2026sentimentc63a31c8}, Analysis of Sentiment Classification of Hotel Reviews Based on Multinomial Naive Bayes \cite{zhirui2020analysis7e9df95c}, Twitter sentiment classification using Naive Bayes based on trainer perception \cite{ibrahim2015twitter8ea1c52b}, Sentiment Classification of Student Opinions on AI Utilization Using Naive Bayes Algorithm \cite{yuniar2025sentimentaf01069c}, Sentiment Analysis of Wordpress Application Satisfaction Level Using K-Nearest Neighbor and Naive Bayes Methods \cite{hamid2025sentimente313680b}, Sentiment Analysis and Social Cognition Engine (SEANCE): An automatic tool for sentiment, social cognition, and social-order analysis \cite{crossley2016sentimente6f85bf3}, A multimodal approach to cross-lingual sentiment analysis with ensemble of transformer and LLM \cite{miah2024multimodal6ed04e56}, Sentiment Analysis of Customer Reviews of Food Delivery Services Using Deep Learning and Explainable Artificial Intelligence: Systematic Review \cite{adak2022sentiment1dd3a5fa}, A Survey on Aspect-Based Sentiment Classification \cite{brauwers2021survey18bf5cde}, Sentiment and position-taking analysis of parliamentary debates: a systematic literature review \cite{abercrombie2020sentimenta45f4c5f}, Quantum-inspired multimodal fusion for video sentiment analysis \cite{li2020quantuminspireddea5ae5a}, Using rhetorical structure in sentiment analysis \cite{hogenboom2015usingbcdd5ee6}, BanglaBook: A Large-scale Bangla Dataset for Sentiment Analysis from Book Reviews \cite{kabir2023banglabook8f8d1681}, MULDASA: Multifactor Lexical Sentiment Analysis of Social-Media Content in Nonstandard Arabic Social Media \cite{alwakid2022muldasa494768c0}, Opinion mining with the SentWordNet lexical resource \cite{ohana2009opinionb1086ed1}, Proceedings of the 8th Workshop on Computational Approaches to Subjectivity, Sentiment and Social Media Analysis \cite{joshi2017proceedingsb60a4ff4}, ConKI: Contrastive Knowledge Injection for Multimodal Sentiment Analysis \cite{yu2023conkif6d19a04}, A Survey of Sentiment Analysis and Sarcasm Detection \cite{yacoub2024surveyec309e9e}, Sentiment Analysis of Text Guided by Semantics and Structure \cite{hogenboom2009sentiment12bb51d1}, AMR-based Network for Aspect-based Sentiment Analysis \cite{ma2023amrbased5cc90650}, Oil Sector and Sentiment Analysis—A Review \cite{santos2023oilf649334f}, Neuro-Symbolic Sentiment Analysis with Dynamic Word Sense Disambiguation \cite{zhang2023neurosymbolic2490f8a5}, A systematic review on automated clinical depression diagnosis \cite{mao2023systematic791ed455}, SENTIMENT ANALYSIS IN SOCIAL MEDIA: HOW DATA SCIENCE IMPACTS PUBLIC OPINION KNOWLEDGE INTEGRATES NATURAL LANGUAGE PROCESSING (NLP) WITH ARTIFICIAL INTELLIGENCE (AI) \cite{alam2025sentiment284f7433}, Deep Lexical Hypothesis: Identifying personality structure in natural language \cite{cutler2022deep59d94f2d}, Is neuro-symbolic AI meeting its promises in natural language processing? A structured review \cite{hamilton2022neurosymbolicf89bd5ec}, Retrieval-Augmented Generation (RAG) in Healthcare: A Comprehensive Review \cite{neha2025retrievalaugmented537affb3}, Tsetlin Machine for Sentiment Analysis and Spam Review Detection in Chinese \cite{zhang2023tsetline116ff6e}, USSA: A Unified Table Filling Scheme for Structured Sentiment Analysis \cite{zhai2023ussa558da815}, Multifaceted Sentiment Detection System (MSDS) to Avoid Dropout in Virtual Learning Environment using Multi-class Classifiers \cite{t2023multifaceted30921198}, Detection of Autism Spectrum Disorder (ASD) Symptoms using LSTM Model \cite{mukherjee2024detection1f9cc3ec}, Revisiting SMS Spam Detection: The Impact of Feature Representation on Classical Machine Learning Models \cite{soysald2026revisiting0c8a8c5a}, Language Modeling on Tabular Data: A Survey of Foundations, Techniques and Evolution \cite{ruan2024languageed4af37b}, Sentiment analysis on Twitter \cite{proscia2017sentiment7ae0491a}, DemoShapley: Valuation of Demonstrations for In-Context Learning \cite{xie2024demoshapley0eb98bc5}, Representation, information extraction, and summarization for automatic multimedia understanding \cite{harrando2022representation75f8b324}, Sentiment analysis of trending tweets using Spark NLP and deep learning: a benchmark study of CNN vs transformer models \cite{benyamani2026sentiment9b767c99}, Sparse-by-Design Cross-Modality Prediction: L0-Gated Representations for Reliable and Efficient Learning \cite{cenacchi2026sparsebydesign83f7d32c}, Shapley values as a generic approach to interpretable feature selection \cite{trotskii2023shapley828133b0}, Position:TrustLLM: Trustworthiness in Large Language Models \cite{yue2024positionfd8c4283}, Table 1: Sentiment classification performance (Accuracy/Macro-F1) on four financial tweet datasets. \cite{paperndtable3751644a}, Figure 10: The ablation experiment for the positive and negative emotion on IMDB datasets. \cite{paperndfigure6b82e845}, Figure 11: Macro F1 (\%) score improvement across ablation experiments for CICIDS 2017 and CTU-IoT-malware datasets with ± SD error bars (5 × 5-fold CV). \cite{paperndfigure49caf9ce}, Develop a Neural Model to Score Bigram of Words Using Bag-of-Words Model for Sentiment Analysis \cite{balamurali2022develop3ea66269}, Sentiment Analysis of Yelp using Advanced V Model \cite{shaout2019sentiment15df16e4}, Analyze IMDb movies by sentiment and topic analysis \cite{ouyang2023analyze6d51f5aa}, Sentiment Analysis of IMDB Movie Reviews: Natural Language Processing for Opinion Mining \cite{paper2026sentiment19dab9d7}, Extraction of Aspects and Opinion Indicators from Yelp Reviews Using Different Methods of Sentiment Analysis \cite{bizel2020extraction580394cd}, Sentiment Analysis of IMDb Movie Reviews \cite{ajmera2022sentiment9a5a9e51}.

\section{Method}
The configured domain experiment workspace executes the prespecified trials.
Each trial writes structured metrics, and the pipeline keeps only metric improvements.

\section{Experiments}
\noindent `baseline-majority-seed0`: metric=0.3381535038932147, decision=keep, status=ok.\\
\noindent `baseline-length-seed1`: metric=0.42281550405704854, decision=keep, status=ok.\\
\noindent `baseline-lexicon-seed2`: metric=0.4807446118530743, decision=keep, status=ok.\\
\noindent `baseline-unigram-seed3`: metric=0.8335965241857798, decision=keep, status=ok.\\
\noindent `ablation-unigram-alpha05-seed1`: metric=0.8403343304493236, decision=keep, status=ok.\\
\noindent `ablation-unigram-minfreq2-seed2`: metric=0.8268594562677809, decision=discard, status=ok.\\
\noindent `ablation-bigram-seed3`: metric=0.8285593220338983, decision=discard, status=ok.\\
\noindent `ablation-bigram-minfreq2-seed4`: metric=0.8285709443427565, decision=discard, status=ok.\\

\section{Results}
The best kept trial was `ablation-unigram-alpha05-seed1` with primary metric 0.8403343304493236.
\title{Research Decision}
\maketitle

Decision: PROCEED

Proceed with trial `ablation-unigram-alpha05-seed1` as the current evidence baseline.

\section{Limitations}
Claims are limited to the registered assets, evaluation units, trials, metrics, and compute budget recorded in this run.

\section{Conclusion}
The workflow now links literature, experiment metrics, and paper text through auditable artifacts.

\section{Revision Notes}
The revision preserves only screened citations and verified experiment metrics. Open reviewer concerns are tracked in `reviews.md`.

<!-- Review summary preserved for audit -->
<!-- \# Reviews

\section{Reviewer A: Methodology}
\noindent The paper is structurally complete and reports the current evidence.\\
\section{Reviewer B: Evidence}
\noindent Claims should remain tied to verified metrics and screened citations.\\
\section{Reviewer C: Venue Readiness}
\noindent Venue readiness remains conditional on exact template and policy checks.\\
 -->
\end{document}
